\documentclass[12pt,a4paper]{article}

\usepackage[utf8]{inputenc}
\usepackage[T1]{fontenc}
\usepackage[margin=2.5cm]{geometry}
\usepackage{hyperref}
\usepackage{microtype}
\usepackage{parskip}
\usepackage{lmodern}

\hypersetup{colorlinks=true, urlcolor=blue!60!black}

\pagestyle{empty}

\begin{document}

\begin{flushright}
Kundai Farai Sachikonye \\
Auf der Lichtung 19, 82178 Puchheim \\
\href{https://github.com/fullscreen-triangle}{github.com/fullscreen-triangle} \\
\texttt{kundai.sachikonye@bitspark.com} \\
(+49) 1626386344 \\[6pt]
16 May 2026
\end{flushright}

\vspace{1em}

University of Maryland Global Campus \\
Adjunct Faculty Hiring --- Artificial Intelligence \\
Grafenwöhr, Germany

\vspace{1.5em}

\textbf{Re: Adjunct Faculty Position --- Artificial Intelligence}

\vspace{1em}

Dear Hiring Committee,

I am applying for the Adjunct Faculty position in Artificial Intelligence. My
subject-matter depth spans the full stack of what an AI curriculum covers ---
from mathematical foundations and model training to agentic system design ---
and I have built original systems at each layer, including a novel operating system
with its own theoretical framework and kernel language.

At the systems level, \href{https://github.com/fullscreen-triangle/buhera}{Buhera}
is a research OS whose architecture rests on three theorems I developed: the
Triple Equivalence (observation, computation, and processing are formally identical
operations), the S-Entropy Uniqueness Theorem (a provably unique continuous embedding
for backward trajectory navigation), and the Circularity Theorem (why content-based
self-preparation cannot produce a non-circular progress measure). The kernel includes
a Proof Validation Engine using Lean~4, a Penultimate State Scheduler that prioritises
processes by distance to goal state rather than by time slice, and a declarative
kernel language (vaHera) in which programs specify final states rather than
instruction sequences. Validated end-to-end: 100\% accuracy on 32 NIST compounds,
$O(\log_3 N)$ retrieval latency, 122,640$\times$ speedup over discrete search at
the embedding boundary. Six peer-review-ready papers are in the repository.
This gives me the ability to teach AI foundations --- including computability,
complexity, and the theoretical limits of what learning systems can and cannot do ---
from first-hand derivation rather than textbook summary.
The Triple Equivalence is not merely a theorem on a slide: the
\href{https://space-observatory-xi.vercel.app/}{Space Observatory} at
\textit{space-observatory-xi.vercel.app} is its empirical proof. GPU fragment
shaders compute astronomical observables --- Rayleigh/Mie scattering, orbital
mechanics, atmospheric optics --- with median relative error $6.9 \times 10^{-5}$
(10/10 benchmarks); and Newton's $G$ is derived from bounded phase-space partition
structure via three independent routes, matching CODATA~2018 at depth $n=8$ and
double-precision machine epsilon at $n=27$ (43/43 benchmarks).
Astronomy without a telescope; $G$ without an experiment.
This is the kind of first-hand result that can ground an AI foundations lecture on
what ``computation'' means at a level that goes beyond ``matrix multiplications on
a GPU.''

For a computer networks or cybersecurity module, the
\href{https://github.com/fullscreen-triangle/pylon}{Pylon} framework provides
an equally concrete teaching example. Theorem~3 of the framework proves a
bijective symplectomorphism between distributed network phase space and ideal
gas phase space --- this is not an analogy, it is a structural identity.
$PV = Nk_BT$ holds exactly for networks: pressure is communication load,
volume is address space, temperature is timing variance.
Theorem~17 (Central Molecule Impossibility) derives from Landauer's principle
$\epsilon = k_B T \ln 2$ that per-packet state tracking requires infinite entropy ---
the formal reason why per-acknowledgment protocols are thermodynamically
incoherent at scale, not merely inefficient.
Security follows from the Second Law: any attack that reduces network entropy
below its thermodynamic minimum is physically forbidden, giving 100\% detection
with 0\% false positives (128-node, $10^6$-packet validation).
Byzantine fault tolerance achieves 51\% (vs PBFT's 33\%) as a theorem, not
as an empirical result.
Throughput $33\times$, jitter $20\times$, recovery $1000\times$ --- all
validated against prediction to within 2.4\%.
Students who understand why a network is an ideal gas understand distributed
systems at a fundamentally different level than students who memorise TCP/IP layers.

At the model training level, \href{https://github.com/fullscreen-triangle/purpose}{Purpose}
formalises the \textbf{Backward Training Principle} --- a formal inversion of
conventional easy-to-hard curriculum learning proved strictly optimal for any domain
admitting a \emph{pure-intent regime}.
The central result: a policy $\epsilon$-competent on the pure-intent regime (scalar
objective, full capability envelope exercised, unambiguous feedback) is
$\epsilon$-competent on every constrained sub-regime by \emph{feasibility inclusion},
not by generalisation.
A Forward Training Collapse Theorem establishes a provably positive
Vapnik--Chervonenkis failure gap for the conventional direction; a Sample Complexity
Theorem gives $(N+1)\times$ sample savings for a cascade of depth~$N$, and
$2^{N-1}\times$ for binary-nested cascades.
The origin was concrete: I derived the theoretical minimum lap time for a Formula
circuit and realised the circuit is the pure-intent regime for all driving sub-tasks ---
road, urban, and parking driving are strict feasibility restrictions of it.
The principle generalises to protein design, scientific query compilation, and any
domain with additive constraints and a natural extremal objective.
This is a result I can teach from first derivation: not ``start easy,'' but
``start at the fully-resolved objective and project downward.''

At the interface and agentic layer,
\href{https://github.com/fullscreen-triangle/zangalewa}{Zangalewa} is the
\textbf{Minimum Sufficient Interceptor} (MSI) sitting at the boundary of the Buhera OS.
The OS presents a blank screen --- no windows, no running applications.
A user types; the MSI intercepts the utterance, extracts semantic coordinates in a
bounded $[0,1]^3$ partition space via a hierarchical $k$-ary specialist cascade
requiring $O(\log_k N)$ small-model invocations, and hands typed coordinates to the
vaHera kernel.
A \textbf{Sufficiency Theorem} proves the interceptor's parameter count is
$O(D \cdot b \cdot \log|\Sigma|)$ --- critically independent of the downstream
substrate's state-space size: an OS-scale substrate needs ${\sim}10^7$--$10^8$
parameters, not a $100\,$B-parameter model.
Session trajectory maintenance via precision-by-difference timing and a proved-complete
four-state focus-arbitration machine ($\mathsf{Blank}/\mathsf{Prompt}/\mathsf{Artifact}/\mathsf{Artifact+Prompt}$)
complete the mediation architecture.
The autonomous code-repair pipeline that AI students typically first encounter ---
exception capture $\to$ LLM diagnosis $\to$ patch $\to$ git-isolated apply $\to$
rollback --- is a secondary domain application of this same architecture.
\href{https://github.com/fullscreen-triangle/four-sided-triangle}{Four-Sided-Triangle}
is a separate 8-stage RAG system demonstrating hierarchical cascade routing at the
document-retrieval level, with a Rust backend providing 10--50$\times$ speedup.

Two further frameworks provide teaching examples for database systems and multi-agent
coordination respectively.
\href{https://github.com/fullscreen-triangle/bloodhound}{Bloodhound} and its formal
theory (Lagrangian Bloodhound Agent, LBA) answer the question ``what is a query?''
from first principles.
The PDSVM (Purpose-Driven Shader Virtual Machine) treats every scientific data
object as a bounded oscillatory system described by three S-entropy coordinates
$(S_k, S_t, S_e) \in [0,1]^3$; every query falls into exactly one of three shader
primitive classes (Completeness Theorem); and retrieval via ternary trie traversal
runs in $O(k)$ time independent of database size $N$ --- 373$\times$ speedup at
$N=10^5$, browser-deployable via WebGPU.
The stable slice $T^*$ is a generative constraint solver, not a statistical predictor:
the forcing conditions of the data (structural constraints that \emph{must} be true)
are recovered without approximation.
This is exactly the kind of result that grounds a database module in something students
cannot reduce to ``run a SQL query'' --- it explains what a query is doing in terms
of information geometry.
For multi-agent systems, LBA's Intelligence Index $I(\mathcal{A})$ classifies
agent behaviour in $O(1)$ --- no simulation required --- and identifies six failure
phenotypes (entropy collapse, entropy explosion, phase desynchronisation, purpose drift,
aperture collapse, coordination deadlock) that account for every observed failure mode
in deployed multi-agent systems. The Common-Cell Convergence Theorem shows that agents
with disjoint internal representations can still converge on the same action cell:
coordination is in outcome space, not representation space.
For students building or evaluating agentic systems, these are operational concepts,
not classroom abstractions.

A parallel body of work --- the \href{https://github.com/fullscreen-triangle/fourth-stomach}{Fourth-Stomach} framework ---
applies the same mathematical primitives to financial markets. A bounded economic agent
is modelled as a receiver $(K, D, \pi, \beta)$ with an irreducible noise floor
$\beta > 0$, and from this single primitive the framework derives, as strict
mathematical theorems rather than assumptions: Arrow's impossibility theorem,
the Grossman--Stiglitz paradox, dual-process cognition (System~1/System~2), and
the Efficient Market Hypothesis as the exact condition $\sum \log(1/q_i) = \infty$.
Market equilibrium exists (proved via Banach contraction, no convexity required);
portfolio weights satisfy a discrete Kirchhoff law on the asset correlation graph;
and the same Composition-Inflation formula $T(n,3) = 3 \cdot 4^{n-1}$ that governs
partition depth in the OS architecture reduces a 500-asset ETF rebalance from
$117\,\mathrm{ms}$ to sub-10\,ns via pre-computed lookup tables --- five orders of
magnitude. Six papers, all validated to machine precision ($\leq 10^{-14}$ relative
error). Arrow's impossibility is usually presented as a given; I can explain
where it comes from.

\textbf{Autonomous systems and backward trajectory completion.}
The \href{https://github.com/fullscreen-triangle/verum}{Verum} framework
provides a third class of teaching example: why prediction-based autonomous
systems fail, and what replaces prediction. The conventional AV pipeline
(perception $\to$ prediction $\to$ planning $\to$ control) fails because
exponential divergence of forward trajectories is a structural property of
chaotic systems, not an engineering deficiency to be corrected with more
data. In partition space the Lyapunov exponent is exactly zero ($\lambda_\mathrm{partition} = 0$);
the system completes trajectories \emph{backward} from target states rather
than predicting forward from current state. The safety guarantee follows
directly: ``when the triple equivalence does not converge, the vehicle stops''
--- convergence failure maps to halt, not to unpredictable behaviour, provably.
The pure-intent regime for all driving is the Formula One circuit: road
driving, urban driving, and parking are strict feasibility restrictions of
the F1 envelope. I validated this on real 2023 Bahrain Grand Prix telemetry
(Max Verstappen, 242 samples per lap, 20-node oscillatory circuit graph):
state reconstruction residual $7.2 \times 10^{-9}$; ICE-Turbo fault
detected 13~laps before the synthetic injection event; optimal lap time
prediction $80.4\,\mathrm{s}$ vs.\ observed $82.5\,\mathrm{s}$
($2.5\,\%$ gap).
For students working in defence and operational environments where GPS
denial, sensor spoofing, and adversarial conditions are real failure modes,
the difference between a system that predicts and a system that completes
is the difference between one that degrades gracefully and one that
fails catastrophically.

UMGC's student population --- military personnel and working professionals --- benefit
most from instruction that connects abstract concepts to real engineering decisions.
I can teach why attention is $O(n^2)$ in sequence length and what that means for
practical deployment. I can explain why RAG exists as a solution to a specific failure
mode in fine-tuning, and when you should fine-tune instead. I can walk through why
autonomous agents fail in ways that simple chatbots do not, using failure modes I have
encountered in systems I built. This is a different kind of instruction from someone
who has studied these topics; it is instruction from someone who has made the mistakes.

I hold a Master's in Computational Biology (Universität Tübingen) and completed doctoral
research in Computational Lipidomics at TUM. I am legally resident in Germany with full
work authorisation and am available immediately. English is native; German is conversational.

\vspace{1.5em}
Yours sincerely,

\vspace{2.5em}
\textbf{Kundai Farai Sachikonye}

\end{document}
